
\section{HiPPO-LegS: Scaled Measures for Timescale Robustness}
\label{sec:theory_legs}

Exposing the tight connection between online function approximation and memory allows us to produce memory mechanisms with better theoretical properties, simply by choosing the measure appropriately.
Although sliding windows are common in signal processing (\cref{sec:rw-ops}), a more intuitive approach for memory should \emph{scale} the window over time to avoid forgetting.

Our novel \textbf{scaled Legendre measure (LegS)} assigns uniform weight to all history $[0, t]$:
$\mu^{(t)} = \frac{1}{t} \mathbb{I}_{[0, t]}$. %
App~\ref{sec:derivations}, Fig.~\ref{fig:measures} compares LegS, LegT, and LagT visually, showing the advantages of the scaled measure.


Simply by specifying the desired measure, specializing the HiPPO framework (Sections~\ref{subsec:hippo-framework}, \ref{sec:discretization})
yields a new memory mechanism (proof in \cref{sec:derivation-legs}).
\begin{theorem}
  \label{thm:legs}
  The continuous- \eqref{eq:scaled-legendre-dynamics}
  and discrete- \eqref{eq:legs-discrete} time dynamics for \textbf{HiPPO-LegS} are:
  \small

  \begin{minipage}{.35\linewidth}
    \begin{align}
      \ddt c(t) &= -\frac{1}{t} A c(t) + \frac{1}{t} B f(t)
      \label{eq:scaled-legendre-dynamics}
      \\
      c_{k+1} &= \left( 1-\frac{A}{k} \right) c_k + \frac{1}{k} B f_k
     \label{eq:legs-discrete}
    \end{align}
  \end{minipage}%
  \hfill
  \begin{minipage}{.52\linewidth}
    \begin{align*}
      A_{nk}
      =
      \begin{cases}
        (2n+1)^{1/2}(2k+1)^{1/2} & \mbox{if } n > k \\
        n+1 & \mbox{if } n = k \\
        0 & \mbox{if } n < k
      \end{cases},
      \qquad
      B_n = (2n+1)^{\frac{1}{2}}
    \end{align*}
  \end{minipage}
\end{theorem}

We show that HiPPO-LegS enjoys favorable
theoretical properties: it is invariant to input timescale, is fast to compute,
and has bounded gradients and approximation error.
All proofs are in \cref{sec:hippo-theory}.

\paragraph{Timescale robustness.}
As the window size of LegS is adaptive, %
projection onto this measure is intuitively robust to timescales.
Formally, the HiPPO-LegS operator is \emph{timescale-equivariant}:
dilating the input $f$ does not change the approximation coefficients.
\begin{proposition}
    \label{prop:timescale}
    For any scalar $\alpha > 0$, if $h(t) = f(\alpha t)$, then
    $\hippo(h)(t) = \hippo(f)(\alpha t)$.
    \\ In other words, if $\gamma : t \mapsto \alpha t$ is any dilation function, then
    $\hippo(f \circ \gamma) = \hippo(f) \circ \gamma$.
\end{proposition}

Informally, this is reflected by HiPPO-LegS having \emph{no timescale hyperparameters};
in particular, the discrete recurrence \eqref{eq:legs-discrete} is invariant to the discretization step size.%
\footnote{\eqref{eq:legs-discrete} uses the Euler method for illustration; HiPPO-LegS is invariant to other discretizations (\cref{sec:discretization-full}).}
By contrast, LegT has a hyperparameter $\theta$ for the window size, and both LegT and LagT have a step size hyperparameter $\Delta t$ in the discrete time case.
This hyperparameter is important in practice; \cref{subsec:low_order_projection} showed that $\Delta t$ relates to the gates of RNNs, which are known to be sensitive to their parameterization~\cite{jozefowicz2015empirical,tallec2018can,gu2020improving}.
We empirically demonstrate the benefits of timescale robustness in \cref{subsec:exp-timescale}.









\paragraph{Computational efficiency.}
In order to compute a single step of the discrete HiPPO update, the main operation is multiplication by the (discretized) square matrix $A$.
More general discretization specifically requires fast multiplication for any matrix of the form $I + \Delta t \cdot A$ and $(I - \Delta t \cdot A)^{-1}$ for arbitrary step sizes $\Delta t$.
Although this is generically a $O(N^2)$ operation, LegS operators use a fixed $A$ matrix with special structure
that turns out to have fast multiplication algorithms for any discretization.%
\footnote{It is known that large families of structured matrices related to orthogonal polynomials are efficient~\cite{de2018two}.}
\begin{proposition}%
    \label{prop:efficiency}
    Under any generalized bilinear transform discretization (cf.\
    \cref{sec:discretization-full}), each step of the HiPPO-LegS recurrence in equation~\eqref{eq:legs-discrete} can be computed in $O(N)$ operations.
\end{proposition}

\cref{subsec:exp-scalability} validates the efficiency of HiPPO layers in practice, where unrolling the discretized versions of \cref{thm:legs} is 10x faster than standard matrix multiplication as done in standard RNNs.

\paragraph{Gradient flow.}
Much effort has been spent to alleviate the \emph{vanishing gradient problem} in RNNs~\cite{pascanu2013difficulty},
where backpropagation-based learning is hindered by gradient magnitudes decaying exponentially in time.
As LegS is designed for memory, it avoids the vanishing gradient issue.
\begin{proposition}%
  \label{prop:gradient-bound}
    For any times $t_0 < t_1$, the gradient norm of HiPPO-LegS operator
    for the output at time $t_1$ with respect to input at time $t_0$ is
    $\left\| \frac{\partial c(t_1)}{\partial f(t_0)} \right\| = \Theta\left( 1/t_1 \right)$.
\end{proposition}


\paragraph{Approximation error bounds.}
The error rate of LegS decreases with the smoothness of the input.

\begin{proposition}%
  \label{prop:approximation_error}
  Let $f \colon \mathbb{R}_+ \to \mathbb{R}$ be a differentiable function, and let
  $g^{(t)} = \proj_t(f)$ be its projection at time $t$ by
  HiPPO-LegS with maximum polynomial degree $N-1$.
  If $f$ is $L$-Lipschitz then $\norm{f_{\leq t} - g^{(t)}} = O(tL/\sqrt{N})$.
  If $f$ has order-$k$ bounded derivatives then
  $\norm{f_{\leq t} - g^{(t)}} = O(t^k N^{-k+1/2})$.
\end{proposition}

